\problem{P25}{直径与偏心弦的分区}{交点计数}
单位圆内画出 $n\ge1$ 条直径，其方向独立且均匀分布于 $[0,\pi)$。再画一条预先固定的弦，它所在直线到圆心的距离为 $d\in(0,1)$。求最终分区数 $R$ 的分布、均值和方差。
\solution 以概率一，直径方向互不相同；它们全部经过圆心，形成 $2n$ 个区域。令固定弦所在直线为 $x=d$。方向角为 $\theta$ 的直径与它相交于圆内，当且仅当
\[
\frac d{|\cos\theta|}<1.
\]
因此每条直径与弦相交的概率为 $p=2\arccos(d)/\pi$。各方向独立，交点数 $K$ 服从 $\operatorname{Bin}(n,p)$。因 $d>0$，弦不经过圆心，不同直径在弦上的交点互不相同，端点相交也为零概率事件。

$K$ 个交点把新弦分成 $K+1$ 段，每段将所在区域一分为二，所以
\[
\boxed{R=2n+K+1.}
\]
于是对 $k=0,\ldots,n$，
\[
\Prob(R=2n+k+1)=\binom nkp^k(1-p)^{n-k},
\]
且
\[
\boxed{\E R=2n+1+np,\qquad \Var R=np(1-p).}
\]
\pitfall 这里原有直径共点，不能使用“每对直线贡献一个不同交点”的通用位置公式。$d=0$ 也被排除，因为此时新弦经过同一个中心，交点不再互异。
